\section*{Working Notes: Points to Revisit}
\begin{itemize}
\item \textbf{Parent and grandparent averages in Table A1.}
Revisit the erroneously labelled average-parent and average-grandparent
entries in the book's technical appendix. Distinguish a correlation with
an average, a regression slope on that average, and an average of
individual-relative correlations. Averaging changes reliability and the
normalization, so the common individual attenuation factor cannot simply
be carried over. Check the practical relevance separately: these average
rows do not appear among the Table 4.5 moments specified for estimation.
\item \textbf{Intuition for nonidentification of measurement error under
genetic assortment.} Explain why independent true environmental variation
and independent recording noise both dilute the observed genetic signal
without affecting mate selection on genetic value. The correlations reveal
only the product $\theta h^2$: less reliable measurement of a more
heritable trait is observationally equivalent to more reliable measurement
of a less heritable trait. Even the observed spouse correlation does not
separate the two in this model.
\item \textbf{The fundamental measurement-error exclusion.}
An imperfect measure of a latent trait is not automatically that trait
plus classical measurement error. The gap between observed education and
an asserted latent social ability may reflect family resources, access
to schooling, preferences, discrimination, and sex-specific opportunities.
These influences may covary with the person's own genetic and environmental
components, relatives' components, and relatives' measurement residuals.
Our common-$\theta$ correction sets these covariances to zero. If they
are nonzero, additional terms enter the observed kinship covariances and
a uniform attenuation factor need not apply. Develop this objection:
an admissible estimate of $\theta$ does not validate the exclusions or
establish that an error-free heritability has been recovered.
\item \textbf{Identification with generation-specific measurement error.}
Revisit identification when attenuation differs across generations rather
than being governed by one common $\theta$. How can the difference between
full-sibling and parent--child correlations then distinguish the lineal
factor $(1+r)/(1+m)$ from different attenuation factors? The same question
arises for the single-grandparent moment relative to the aunt/uncle moment.
Specify whether reliability varies by historical generation or by the
number of generations separating a pair, derive the resulting pair-level
attenuation factors, and establish what additional restrictions or moments
would separate measurement error from the effects of $r$ versus $m$.
\item \textbf{Measurement correction added to A1--A3, but parameter recovery unclear.}
The book inserts $\theta$ into the moment equations, so the A3 intercept
identifies $\exp(a)=\theta h^2$, not latent heritability alone. Yet the
surrounding text still describes estimated heritability and cites the PNAS
SI Table S2. Check whether the plotted/reported quantity is $\theta h^2$
or a separately recovered $h^2$. Under phenotypic assortment, recovery
from A3 is possible using $m=2b-1$, $r=(1+m)\exp(c)-1$, $h^2=m/r$,
and $\theta=\exp(a)/h^2$, subject to admissibility; retaining the same
regression is not itself an identification error. Distinguish this
passage from Chapter 5's use of Table A2's spouse and affine-relative
moments for Denmark and Sweden. Under the stated common-attenuation
formulas, the sibling-in-law/sibling ratio identifies $r$, while the
spouse correlation identifies $\theta r$, providing an additional route
to separate reliability and heritability. Audit that estimation and
Table 5.4 separately rather than assuming the PNAS procedure applies.
\item \textbf{Reconcile Figure A1, PNAS Figure 1, and PNAS SI Table S2.}
The book's Figure A1 (p.~298) is the counterpart of PNAS Figure~1,
but several plotted coordinates differ. Approximate graphical readings
of the heritability axis put higher education 1860--1919 near 0.58 in
PNAS and 0.33 in the book, and occupational status 1860--1919 near
0.66 and 0.53, respectively. PNAS Figure~1 also disagrees with its own
SI Table S2: occupational status 1780--1859 is plotted near 0.94,
whereas the table reports 0.722. These are graphical readings, not
recovered numerical estimates. Establish the source data, estimation
version, and parameter interpretation for each display; do not infer
from the changed points alone that measurement error or the structural
constraint was newly handled correctly. For an independent extension,
the book's Table 5.1 (p.~67) reports observed education correlations for
spouses and several affine relatives in Denmark and Sweden. Keep these
distinct from the English education cohorts and from the latent
correlations inferred in Table 5.2.
\item \textbf{An apparent impasse between assortment interpretations.}
With phenotypic assortment and free measurement attenuation, the lineal
contrast and $m=h^2r$ allow separate recovery of $h^2$ and $\theta$.
However, our unweighted A3 fits to Table 4.5 yield inadmissible structural
recoveries for four of eight outcomes: $r>1$ for company directorship,
and $h^2>1$ for occupational status 1780--1859, higher education
1780--1859, and literacy. Alternatively, if the intended assumption is
genotypic assortment, the moment formulas must change: the separate
lineal adjustment disappears ($c=0$), and these kinship correlations
identify only $m$ and $\theta h^2$. Intergenerational observations can
improve estimation and test the common decay pattern, but do not add
the identifying contrast needed to separate $h^2$ from $\theta$.
Thus one cannot rescue the interpretation merely by switching assortment
assumptions while retaining separately identified heritability. Is this
an impasse for the interpretation of the reported estimates? Distinguish
that concern from a formal rejection: bounded phenotypic estimation
remains possible, the results depend on weighting and sampling uncertainty,
and externally justified information about reliability or heritability
could resolve the genotypic model's nonidentification.
\item \textbf{Claims about regression to the mean.}
Otto et al. explicitly discuss linear regression toward the mean in
cultural models, including the possibility of reproducing the genetic
model's rate. Clark is therefore wrong to present linear regression
toward the mean as a distinct prediction of genetic transmission.
Distinguish the shape of the regression relationship from its numerical
rate, and distinguish a fixed prediction from a fitted parameter: in the
genetic benchmark at issue the coefficient is 0.5, whereas in the cultural
model it is a free parameter, which can also take that value. This is a
difference in predictive restriction, not a demonstration that cultural
transmission cannot generate the same pattern. Supply the explicit
passages from Otto et al. and Clark, and specify the assumptions behind
the genetic 0.5 benchmark rather than treating it as universal.
\item \textbf{Key points to hammer home.}
There is genuine value in Clark's analyses, and the review should
acknowledge that explicitly. The central concern is not that the analyses
are without merit, but whether the conclusions he draws follow logically
from them. Preserve this distinction when assessing both contributions
and limitations.
Often Clark's conclusions do not follow from the evidence he presents.
This does not establish that the conclusions are wrong; it means they
have not been established nearly as credibly as he suggests. Develop
this distinction consistently throughout the review. Working analogy
to consider: his apparent resistance to alternative explanations recalls,
in the opposite direction, the stubbornness of Lysenkoite social
scientists who instinctively reject evidence allowing any role for
genetics. The target is asymmetric standards of evidence, not an
equivalence of substantive views. He also sometimes exaggerates how
strongly his tests discriminate between genetic and cultural models;
symmetry and regression toward the mean are key examples to document.
Examine his claims that school reforms yield no benefits, distinguishing
evidence of no benefit from failure to establish a benefit. Track shifts
between policy criteria: does a reform help anyone or improve average
outcomes, does it help the poor, does it reduce inequality, and does it
increase or reduce mobility? These are distinct questions, and a failure
on one criterion does not establish failure on the others. Collect
specific passages showing where the criterion changes within an argument.
\item \textbf{Selective discussion of assumptions in twin and adoption studies.}
Clark's discussion focuses on one assumption whose violation can bias
estimated $h^2$ downward, but these designs depend on other assumptions
as well. Identify the assumption he emphasizes and examine the others,
including violations that could bias estimates upward or have an
ambiguous effect. Distinguish twin from adoption designs and assess the
direction of each bias under explicit conditions. A possible source of
downward bias alone does not establish that the estimates are lower
bounds on heritability.
\item \textbf{Lower resemblance when raised apart and claims about adoptees.}
Add a discussion of the evidence for lower resemblance among relatives
raised apart than among those raised together, and what this comparison
implies for Clark's interpretation of environmental influences. His
conclusions about adoptees appear much stronger than the evidence
warrants. Identify the precise claims and supporting studies, report
the relevant comparisons and uncertainty, and distinguish what those
studies establish from broader claims about the absence of effects of
rearing environments.
\item \textbf{Within-MZ-twin estimates are not automatically causal.}
Acknowledge Clark's valid point that comparisons within monozygotic
twin pairs do not by themselves identify causal effects. Controlling
for shared genes and family background does not eliminate nonshared
confounding or measurement error. Also note that he does not cite the
prior literature documenting these limitations, either within economics
or outside it. Supply the relevant references and check his discussion
and bibliography when developing this point: the methodological concern
is important, but it is not new.
\item \textbf{Implications for education spending and social outcomes.}
Discuss the claims: ``First implication\ldots\ is that modern societies
spend too much on education.'' ``What drives much of this expenditure
is a belief that boosting years of education by less advantaged students
increases social mobility rates.'' ``Second implication\ldots\ social
outcomes will change only when the genetic endowment of the society
changes.'' This implies that maximum genetics explains 60\%.
\item \textbf{The covariance exclusions matter beyond genetic nurture.}
Revisit Clark's rejection of genetic nurture on p.~86, but do not frame
the objection as depending on this particular causal mechanism. Any
relevant covariance that the fitted model incorrectly sets to zero can
bias the recovered genetic, assortment and attenuation parameters.
This includes environmental covariance between relatives and
cross-person genetic--environmental covariance beyond the terms already
implied by the stated assortment model. Establishing or rejecting a
specific ``nature of nurture'' pathway does not establish the required
zero-covariance restrictions. Write out the omitted covariance terms and
show how they enter the fitted moments; do not presume that every
violation biases every parameter, or that the direction is always the
same. The issue is the validity of the identifying exclusions, not the
label attached to the source of resemblance.

\item \textbf{Binary phenotypes and the scale of the model.}
Examine whether the linear phenotype, kinship-correlation and common
measurement-error formulas apply to binary outcomes such as higher
education, literacy and company directorship. Distinguish correlations
and heritability on the observed binary scale from those of a latent
liability. Investigate whether prevalence, thresholds, and differences
across cohorts or sexes alter the predicted moments, attenuation and
identification. In a threshold/liability model, the observed Pearson
correlation is generally a nonlinear function of the latent correlation
and the two marginal prevalences. Check explicitly whether one common
attenuation factor $\theta$ can represent that transformation across
kinship types; generally it need not. Compare a properly specified
threshold model with the direct fit to binary Pearson correlations, and
check whether Clark supplies such a link anywhere in the book or PNAS
materials. Work out the implications before treating a binary outcome
as a continuous phenotype with classical measurement error.
The potentially serious problem is that thresholding changes the pattern
of correlations across kinship types, not merely their level. Even if
Fisher's model fits the latent correlations perfectly, its geometric
decay restrictions generally fail for the observed binary Pearson
correlations; fitting the original formulas can then distort inferred
heritability and assortative mating. Demonstrate this using jointly
normal liabilities whose population correlations satisfy the model
exactly. Threshold them and compare fits to population tetrachoric
correlations (the binary case of polychoric correlations), which recover
the latent correlations under these assumptions, and to Pearson
correlations. At 50\% prevalence in both groups, the exact mapping is
$r_{\mathrm{binary}}=(2/\pi)\arcsin(r_{\mathrm{latent}})$: latent correlations
of $0.8$ and $0.4$ become approximately $0.590$ and $0.262$, so even their
ratio changes. Use several kinship types to show why neither a common
attenuation factor nor refitting a single persistence rate generally
restores the exact fit. The problem is especially consequential for the
$\theta$ extension: classical measurement error motivates a common
multiplicative attenuation, whereas thresholding generally implies an
effective $\theta_k=r_{\mathrm{binary},k}/r_{\mathrm{latent},k}$ that varies
across kinship types and marginal prevalences. In the example above,
these ratios are approximately $0.738$ and $0.655$. On the log scale,
$\log\theta_k$ is therefore not a common intercept shift: it can alter
the fitted kinship-decay slope and lineal--collateral contrast used to
recover the structural parameters. Allowing one free $\theta$ does not
generally repair this misspecification; allowing a separate unrestricted
$\theta_k$ for every kinship type would instead remove the cross-kinship
restrictions that make the model informative. A threshold model should
derive these transformations from prevalences and latent correlations.
This population example isolates measurement-scale
misspecification from sampling error; quantify its practical importance
at the prevalences and correlations relevant to Clark's data.
\item \textbf{Fit competing cultural and genetic transmission models.}
Assess the absence of a serious comparison with joint cultural--genetic
transmission or a purely cultural model. A good fit of the restricted
genetic model, even if established, does not identify that mechanism.
Specify and compare the alternatives under transparent assumptions and
the same data and objective, distinguishing predictive restrictions
from extra flexibility. For the regression-to-the-mean point above,
\citet[pp.~8--9, eq.~8]{OttoEtAl1994} write
\[
 B_o=\beta_I(B_m+B_f)+\delta S_C,
\]
where $S_C$ is cultural transmission error, not schooling. The equation
resembles genetic transmission, but $\beta_I$ is not fixed at one half.
They explicitly describe regression towards the mean under cultural
transmission; this pattern alone cannot select the genetic explanation.
\item \textbf{Disconnect between behavioural genetic estimates and policy conclusions.}
Develop the distinction between variance-component estimates and the
causal effects of feasible interventions. Fixed genes do not imply
fixed effects, and a statistical direct genetic effect need not operate
through an unmodifiable physiological pathway. Use the eyeglasses
example and relate the point to \citet{Goldberger1978} and
\citet{Jencks1980}. Give explicit credit for the provocative questions,
valuable dataset and interesting patterns while explaining why the
model-selection and policy conclusions go beyond what those results
establish. Engage with the empirical evidence rather than personal labels.

\item \textbf{Assortment on ``something deeper'' requires a model.}
Address Clark's argument that people assort on something deeper than
the observed phenotype. Specify what that latent characteristic is,
how it relates to observed outcomes and their genetic and environmental
components, and how it governs partner selection. Derive the spouse
covariances and kinship moment conditions under that matching rule,
including the measurement assumptions needed to connect latent and
observed outcomes. An appeal to a deeper trait is not sufficient to
retain the original moment formulas or parameter interpretations.
If this is simply phenotypic assortment on a latent trait measured with
classical error, state and justify those restrictions explicitly; if it
is a different assortment mechanism, establish which moments and
identification results change.

\item \textbf{Selective treatment of bias and lower within-family GWAS estimates.}
Examine whether Clark acknowledges and adequately addresses the much
lower estimates reported in within-family GWAS. His emphasis on biases
that could depress other estimates should be assessed alongside evidence
that controlling for family background reduces population GWAS effect
estimates or polygenic-score associations. Identify the relevant studies,
report the comparable within-family and population estimates with their
uncertainty, and check his discussion before asserting an omission.
Distinguish SNP effects, score associations or prediction, and SNP-based
heritability from total narrow-sense heritability; these are not
interchangeable quantities. Assess explanations involving indirect
genetic effects, population structure, assortment and measurement error
under explicit assumptions, rather than treating the lower estimates
as either automatically decisive or automatically dismissible.

\item \textbf{Use within-family PGIs in spousal analyses.}
Clarify that analyses intended to interpret spousal PGI resemblance as
assortment on direct genetic effects should use PGIs constructed with
within-family GWAS effect estimates. Population-GWAS weights can also
reflect population structure and indirect genetic associations, so
spousal correlations in those scores need not isolate the intended
genetic component. Compare the results using population and within-family
weights, allowing for differences in precision and predictive accuracy.
Within-family weighting is not a complete solution: establish what the
score measures, which biases remain, and how its spousal correlation
relates to the model's latent additive-genetic correlation $m$.

\item \textbf{Misspecification need not produce poor fit.}
Make explicit that a misspecified model can fit the observed moments
well: fitted parameters may absorb omitted covariance terms, and
different mechanisms may generate similar patterns of resemblance.
Consequently, good fit does not validate the identifying restrictions
or the genetic interpretation of the estimates. Keep this logical point
separate from the empirical concern that fit may itself be poor once
the model's restrictions and admissible parameter bounds are imposed.
Assess both, but do not make the identification criticism depend on
establishing poor fit or a formal statistical rejection.

\item \textbf{Residual stratification despite ancestry restrictions and quality control.}
Document the substantial GWAS evidence that population stratification
can bias estimates even after standard quality control and restriction
to an apparently homogeneous ancestry group. Such restrictions do not
establish the absence of fine-scale population structure or its
association with social and environmental differences. Assemble concrete
examples, specifying the phenotype, ancestry definition, controls and
evidence used to diagnose residual bias. Explain the implications for
interpreting population-GWAS effects and spousal PGI correlations as
purely genetic quantities. Distinguish stratification from indirect
genetic effects and assortative mating; a within-family versus population
difference alone need not identify which mechanism accounts for it.

\item \textbf{The Collado--Clark debate and the identification disputes of the 1970s.}
Explore the parallels with the debates involving Jencks, the Birmingham
group, and Cloninger and colleagues. A recurring problem is identification
of distinct transmission mechanisms from a small set of kinship moments:
similar fits can conceal very different causal models and parameter
interpretations. Check the historical positions and the precise points
of disagreement before drawing the comparison. The irony is how much
progress has occurred in data, molecular measurement and research designs,
while this basic identification problem can reappear when inference rests
on a limited set of resemblance patterns. Explain which newer evidence
can actually distinguish the competing mechanisms, rather than treating
technical progress itself as a solution to identification.

\item \textbf{Goldberger on adding environments to Fisher's scheme.}
Use Goldberger's pointed observation, preserving its specific target:
``The Birmingham models gain empirical validity by tacking on environmental
correlations to Fisher's scheme'' \citep[p.~36]{Goldberger1978}.
Discuss the accompanying coherence problem: adding a shared-environment
term to sibling resemblance requires an account of the process generating
it and its implications for other kinships. Goldberger asks whether shared
sibling environmental variance should reappear for uncles and cousins,
what equates parent--child and sibling environmental shares in another
specification, and why parental transmission of both environments and genes
does not generate gene--environment covariance. Follow his Appendix to
spell out the restrictions for a coherent extension. Do not turn this into
a claim that every shared sibling influence must persist across generations;
the implications depend on how the environment is generated and transmitted.

\item \textbf{Modern evidence on dominance does not establish zero shared environment.}
Section~\ref{sec:genomic-evidence} now develops the distinction between
dominance and gene--environment covariance. Extend its outcome-specific
evidence where appropriate. Identify the relevant findings,
state exactly which claim about dominance they support, and delimit
the outcomes and designs to which they apply. This does not vindicate
the separate claim that shared-environment variance $c^2$ is zero for
all outcomes in adulthood. Evaluate that claim against outcome-specific
evidence and uncertainty, keeping dominance and shared-environment
components conceptually distinct. Support for one modelling assumption
cannot stand in for evidence for the other.
% User's dictated continuation after "or" is still pending; do not infer it.
\item \textbf{Read Alex Young's commentary on Clark.}
Read Young's initial reaction and extended July 2023 commentary on
Clark's PNAS paper, including his discussion of alternative transmission
models, assortative mating and identification. Compare the arguments
with our review before deciding whether to cite or address them.
The posts are dated 26 June and 4 July 2023; original links are retained
in the saved source notes.
Copies and retrieval notes are saved in the project reference library.

\item \textbf{Add analogy to equilibrium restriction from economics.}
Explain why $m=h^2r$ is a consequence of the maintained phenotypic-assortment
model, not an optional additional assumption. Select an economic analogy
that distinguishes primitive assumptions from their implied restrictions,
and distinguishes estimating an unrestricted model to test the restriction
from interpreting its fit as the fit of the maintained model.


\par\medskip
\textit{Preferred option: the consumption Euler equation.}
Consider a consumer with logarithmic utility who can transfer resources
between two periods at gross return $R$, under certainty and with an
interior optimum. Utility maximisation implies $c_2/c_1=\beta R$, where
$\beta$ is the discount factor. Consumption growth, patience and the
return cannot be specified independently while preserving that model.
The Euler equation is a consequence of the optimisation problem, not
an optional additional assumption. The assortment restriction plays an
analogous role: freeing $m$, $r$ and $h^2$ requires changing the maintained
model, rather than simply declining to impose one of its consequences.

\par\medskip
\textit{Preferred option: homogeneity of consumer demand.}
In the standard consumer model, doubling all prices and nominal income
leaves the budget set unchanged. With unchanged preferences and a unique
optimum, it therefore leaves demand unchanged. Homogeneity of demand is
a consequence of the model, not an additional assumption to be accepted
or discarded independently. A fitted demand equation that responds to
this common rescaling may be useful descriptively, but it cannot
represent the maintained choice model. Likewise, the phenotypic-assortment
model does not permit independent choices of $m$, $r$ and $h^2$: its
assumptions require $m=h^2r$.

\item \textbf{Cite the Georgia land lottery paper.}
Cite \citet{BleakleyFerrie2016}, ``Shocking Behavior: Random Wealth in
Antebellum Georgia and Human Capital Across Generations,'' on the 1832
Cherokee Land Lottery, when discussing exogenous wealth shocks and
intergenerational outcomes.

\item \textbf{Regression towards the mean does not identify the transmission mechanism.}
Focus the criticism of Chapter~14 on the inference from an approximately
linear conditional mean to a claim about its causal mechanism.
The figures do show substantial regression towards the mean; do not
base the critique on denying that descriptive finding. A nearly flat
child-on-parent relationship, as for longevity, means especially strong
regression. The dashed lines in the empirical panels are fitted lines,
not the identity line.

Show explicitly that a purely cultural model can generate the same
linear, symmetric pattern: if
$C_o=\gamma(C_f+C_m)/2+U$, with an independent mean-zero social
innovation $U$, then the conditional mean given the parental average
is linear. Under suitable Gaussian assumptions and innovation variance,
this model can reproduce the one-generation joint distribution of
parental and offspring phenotypes in our additive Fisher specialisation.
This is a counterexample to identification from these plots, not a claim
that every extended-pedigree model is equivalent. Assess Clark's
conclusion that the pattern confirms the likely dominance of direct
genetic transmission (p.~216) against a specified cultural alternative,
rather than equating cultural transmission generally with the particular
poverty/privilege-trap pattern sketched on p.~211.

Supporting qualifications: genetic additivity alone does not guarantee
an exactly linear conditional mean; linearity, contraction and symmetry
about a common population mean require distinct assumptions; the high
reported $R^2$ values for grouped means are not individual-level
explanatory power or formal tests against nonlinearities; and symmetry
of log deviations is not equality of ordinary percentage changes.
Keep these qualifications subordinate to the identification criticism.
The source-checked quotations, graphical slope estimates, identity-line
replots and mathematical counterexamples are saved in the
Chapter~14 audit dated 23 September 2026 in the project notes.
% Source: ../notes/chapter14_regression_to_mean_audit_2026-09-23.md

\item \textbf{Chapter 18: equilibrium variance and a numerical correction.}
With fixed segregation variance $S$, the chapter's equation (6) implies
$V(m)=2S/(1-m)$, hence $V(0.6)/V(0)=2.5$, not the fivefold increase
stated on p.~274; the fivefold result is correct for the correlation
of $0.8$ used in \citet[pp.~3--4]{ClarkCummins2022}.
The discrepancy is consistent with an uncorrected carryover from the
working paper, although the revision history has not been verified.
For the final review, emphasize that this is a standard additive-genetic
equilibrium result under a constant-segregation-variance approximation,
whereas interpreting latent social ability as genetic, relating genetic
to phenotypic assortment ($m=h^2r$ under phenotypic homogamy, with
potentially changing $h^2$), and attributing the Industrial Revolution
to the enlarged upper tail require separate justification.
The increase in equilibrium genetic variance under positive assortative
mating is a classical result of quantitative genetics, originating in
Fisher's framework and developed more explicitly in subsequent theoretical
work. Clark's contribution here lies in its proposed historical application,
whose empirical premises and causal interpretation require separate
evaluation. Trace the result through Fisher, subsequent formal treatments
such as Nagylaki's, and textbook accounts, distinguishing Fisher's explicit
results from later clarifications rather than attributing its origin to
a modern exposition.
% Modern reference: Yengo and Visscher (2018), Theoretical Population Biology
% 124:51--60, doi:10.1016/j.tpb.2018.09.002, equations (3)--(5).

\item \textbf{Cousin marriage does not establish weak genetic assortment.}
Examine Chapter~18's inference from cousin marriage to a low equilibrium
spousal genetic correlation (pp.~278--281). Random choice of a cousin
with respect to ability or status is an assumption, not an implication
of cousin marriage; selection within the permitted kinship pool could
change the result. More fundamentally, establish whether the cousin
correlation $[(1+m)/2]^3$ remains valid under recurrent consanguineous
mating before accepting the fixed-point calculation
$m=[(1+m)/2]^3$: repeated related mating creates additional shared
ancestors and paths of genetic relationship that require an explicit
pedigree or population-genetic treatment. Treat this as a derivation to
check, not as a completed demonstration of the corrected equilibrium.
Distinguish genetic relatedness through shared ancestry from correlation
in additive genetic values for a particular trait. The calculation alone
does not establish weaker genetic assortment in the historical societies
concerned, nor their comparative supply of exceptional innovators.

\end{itemize}
